\section{Related Work}
\label{sec:related}

\subsection{Hierarchical and Prototype-Based Coordination}
The comparisons below distinguish the inherited H-SFP interface from the distributional analysis and intermittent-summary mechanisms studied in E-HSFP.
Federated optimization addresses heterogeneous data and local computation through proximal regularization, normalized updates, and control variates~\cite{li2020federated,10.5555/3495724.3496362,pmlr-v119-karimireddy20a}. Hierarchical FL introduces intermediate edge aggregation~\cite{liu2020client,9488756,rana2022hierarchical}, while SplitFed~\cite{thapa2022splitfed} and HSFL~\cite{quan2023hiersfl} combine partitioned execution with federated coordination. Recent work~\cite{10980018} analyzes hierarchical split training and jointly optimizes model splitting and aggregation. These methods exchange activations or models; to our knowledge, none of them trains upper segments exclusively on synthetic features derived from transmitted statistics.

FedProto~\cite{tan2022fedproto} coordinates local representations through class prototypes. FedTGP~\cite{zhang2024fedtgp} learns server-side global prototypes
through an adaptive-margin contrastive objective. FedProto and FedTGP use prototypes to regularize or classify local representations, whereas in \ehsfp{} the prototypes are the only training inputs available to the subsequent model segments, and the learned scorer weights the sources from which those inputs are synthesized. FedDF~\cite{lin2020ensemble} and FedGen~\cite{zhu2021data} transfer knowledge through ensemble distillation and learned synthesis, respectively. \hsfp{}~\cite{tran2026h} uses class-wise moments to generate training inputs for locally optimized edge and cloud segments. \ehsfp{} retains this interface and studies how incomplete and outdated statistics change the distributions.

\subsection{Statistical Replay, Aggregation, and Consistency}
PASS~\cite{9578909} augments stored prototypes with Gaussian noise, while SLCA~\cite{zhang2023slca} samples from class-wise means and covariances for classifier alignment in continual learning. Statistical replay and covariance-based synthesis are therefore established mechanisms; we analyze the distributions they induce when the underlying summaries are incomplete or outdated. FedFA~\cite{zhou2023fedfa} uses local and federation-wide feature
statistics to construct Gaussian perturbations for feature
augmentation. Its extension introduces alignment based on
higher-order feature statistics~\cite{zhou2024featureaugmentation}.
These methods address feature-distribution differences across
clients. \ehsfp{} examines the synthetic training distributions
produced by incomplete and outdated summaries at successive
model boundaries. \ehsfp{} applies timestamped replay to intermittent upstream inputs, combining received and stored moments before aggregation and generation. Its optional classification covariance pools within-class scatter into a full matrix shared across classes; SLCA uses class-specific covariances. The shared-covariance option requires additional scatter statistics.
Adaptive aggregation includes personalized weighting~\cite{zhang2020personalized} and FedLAW~\cite{pmlr-v202-li23s}, which learns model-aggregation weights and a global shrinking factor using server-side proxy data. \ehsfp{} scores statistical sources using metadata and trains the scorer with heuristic bootstrap targets, so the scores serve as reliability proxies.
MOON~\cite{9578660} contrasts representations of the same example under current, global, and previous local models. R-Drop~\cite{liang2021rdrop} aligns predictions from two dropout realizations of one example through bidirectional KL divergence. PRC compares outputs on independent input-side and memory-based synthetic draws. For a fixed deterministic output map, its expected squared penalty includes both output-mean discrepancies and output-variance terms, as detailed in Section~\ref{sec:regularization}. Source masking separately changes the statistical inputs used for synthesis.

\subsection{Asynchronous and Event-Driven Learning}
Asynchronous Federated Optimization~\cite{xie2019asynchronous}, FedBuff~\cite{nguyen2022federated}, and general analyses~\cite{fraboni2023general} study delayed model updates. These analyses treat delay as a perturbation of the parameter or gradient that is applied. In \ehsfp{}, packet age instead perturbs the input distribution of the next model segment. Section~\ref{sec:formulation} quantifies this perturbation as a Wasserstein drift, and Section~\ref{sec:convergence} converts it into a gradient bias while retaining variation in the fresh-prototype reference objectives.
FedLess~\cite{grafberger2021fedless} and LIFL~\cite{qi2024lifl} provide serverless FL platforms that motivate separating execution from persistent learning state. \ehsfp{} studies this separation in round-driven simulations of delays, masking, and execution disruptions; deployment costs are not measured (Section~\ref{sec:discussion}).
